\subsection{Lagrangian products}
\label{subsec:preliminaries-lagrangian-products}

Write
\[
  \mathbb R^4=\mathbb R_q^2\oplus\mathbb R_p^2
  =\{(q_1,q_2,p_1,p_2)\}.
\]
\begin{definition}[Lagrangian product]
\label{def:preliminaries-lagrangian-product}
For planar convex bodies \(K_q,K_p\subset\mathbb R^2\), their Lagrangian
product is
\begin{equation}\label{eq:preliminaries-lagrangian-product}
  K_q\times_L K_p
  =\{(q_1,q_2,p_1,p_2):(q_1,q_2)\in K_q,\ (p_1,p_2)\in K_p\}.
\end{equation}
Thus \(K_q\) occupies the Lagrangian \(q\)-plane and \(K_p\) occupies the
Lagrangian \(p\)-plane.  This is different from taking a product of planar
domains in the symplectic coordinate planes \((q_1,p_1)\) and \((q_2,p_2)\).
\end{definition}

The two decompositions pair the same four coordinates in different ways:
\[
  \underbrace{(q_1,q_2)}_{K_q}\ \big|\
  \underbrace{(p_1,p_2)}_{K_p}
  \qquad\text{versus}\qquad
  \underbrace{(q_1,p_1)}_{\text{first symplectic plane}}\ \big|\
  \underbrace{(q_2,p_2)}_{\text{second symplectic plane}}.
\]
In block notation \(J_0(u,v)=(-v,u)\).  A normalized facet row
\(a_i=(u,0)\) from the \(q\)-factor therefore has \(J_0a_i=(0,u)\) in the
\(p\)-plane, and a row \(a_j=(0,v)\) from the \(p\)-factor has
\(J_0a_j=(-v,0)\) in the \(q\)-plane. Section~\ref{sec:generalized-reeb-orbits-polytopes}
will identify the corresponding Reeb velocity as \(2J_0a_i\). Thus the
characteristic dynamics exchange the two factors; they do not decouple into
independent planar dynamics.

The central examples of the thesis are Lagrangian products of polygons; their
lifted facet rows and the resulting restricted word enumeration are introduced
with the Haim--Kislev formula, where they are first used.
